% =====================================================================
% Materials and Methods -- revision 2026-10-10
% Edited from the authors' 2026-10-10 draft: duplicates removed, timing
% conflict resolved from the procedure log, US spelling, post-closure
% calibration added (Section 2.9). Red text = information still needed.
% =====================================================================
\section{Materials and Methods}

\subsection{Patient}
\label{sec:study_population}
The patient was enrolled in a prospective modeling study of children scheduled for transcatheter VSD closure at the National Center for Women and Children's Health, RSAB Harapan Kita, Jakarta, Indonesia (December 2025 to June 2026). Children older than three months with an echocardiographically confirmed VSD of at least 2~mm were eligible; children with an additional severe cardiovascular abnormality or previous cardiac intervention were excluded. P1 was the first with a complete catheterization and echocardiographic record both before and after closure, and is the subject of this study. P1 was a 3-year-old girl (age 3.17~years, weight 13.4~kg, height 95~cm) with a perimembranous VSD measuring 5.4~mm on the left-ventricular side and 3.7~mm on the right-ventricular side. The shunt was left-to-right and restrictive: the peak LV--RV systolic gradient was 69~mmHg, mean pulmonary arterial pressure was 15~mmHg and $Q_p/Q_s$ was 1.19 by the Fick method. The study was approved by the Institutional Review Board of RSAB Harapan Kita (DP.04.03/D.XXII/4749/2025). Written informed consent was obtained from the patient's legal guardian, and the data were de-identified.

\subsection{Clinical data}
\label{sec:data}
Clinical data were obtained from the catheterization protocol forms and the medical record (Table~\ref{tab:patient_data}). The defect diameter was measured on echocardiography at both ventricular sides; the right-ventricular side, which is the narrower, restrictive orifice, was used as the model input, taken as the mean of the recorded range of 3.63--3.70~mm ($D_{\mathrm{VSD}} = 3.665$~mm). Heart rate was recorded separately for each state (119~bpm before and 108~bpm after closure) and used as a model input in that state. Where the form records three consecutive readings, the median was used.

Pressures were recorded with a hydrophilic angiographic catheter (APT Medical, Shenzhen, China). Mean arterial pressure was computed from the right femoral artery (RFA) catheter readings with the formula of the catheterization protocol, $\overline{P} = (P_{\mathrm{sys}} + 2P_{\mathrm{dia}})/3$; the mean arterial pressure entered on the form is computed from the cuff readings and was not used. Pulmonary flow $Q_p = 4.087$~L/min and $Q_p/Q_s = 1.194$ were measured by the Fick method, with blood oxygen saturations measured on an AVOXimeter 1000E whole-blood co-oximeter (Accriva Diagnostics, San Diego, CA, USA). Systemic flow and shunt flow were derived as $Q_s = Q_p/(Q_p/Q_s) = 3.423$~L/min and $Q_{\mathrm{shunt}} = Q_p - Q_s = 0.664$~L/min, and systemic vascular resistance as $(\overline{P} - P_{\mathrm{RA}})/Q_s$. Pulmonary vascular resistance could not be derived because left atrial pressure was not measured.

Post-closure pressures were recorded by catheter 9--26~min after device release, in the same session and under the same anesthesia as the pre-closure measurements. Post-closure systemic flow, $Q_p/Q_s$ and ventricular volumes were taken from the post-closure echocardiography in the medical record \textcolor{red}{[timing of this echocardiogram]}: left-ventricular volumes by the biplane method and right-ventricular volumes by three-dimensional echocardiography. The protocol form also records left-ventricular volumes measured with the device deployed but not yet released (32.0/23.6~mL); this alternative record is evaluated in Section~\ref{sec:results_post}.

\begin{table}[!htbp]
\centering\small
\setlength{\tabcolsep}{3.5pt}
\caption{Clinical record of P1 used in this study. Values marked with an asterisk are the targets of the pre-closure calibration; all other values are compared with the model only. No post-closure value was used in the pre-closure calibration.}
\label{tab:patient_data}
\resizebox{\linewidth}{!}{\begin{tabular}{llrr}
\toprule
\textbf{Variable} & \textbf{Method} & \textbf{Pre-closure} & \textbf{Post-closure} \\
\midrule
\multicolumn{4}{l}{\emph{Model inputs}} \\
Age (yr) & medical record & 3.17 & 3.17 \\
Weight (kg) & measured & 13.4 & 13.4 \\
Heart rate (bpm) & ECG monitor & 119 & 108 \\
VSD diameter, RV side (mm) & echocardiography & 3.665 & -- \\
\midrule
\multicolumn{4}{l}{\emph{Pressures (mmHg)}} \\
Systolic / diastolic pressure, RFA & catheter & 100$^{*}$ / 57$^{*}$ & 89 / 68 \\
Mean arterial pressure, RFA & $(P_{\mathrm{sys}} + 2P_{\mathrm{dia}})/3$ & 71.33$^{*}$ & 75.00 \\
Systolic / diastolic pressure, DAo & catheter & 91 / 57 & 104 / 72 \\
Systolic / diastolic pressure, cuff & oscillometric & 119 / 83 & 87 / 55 \\
PA pressure, sys / dia / mean & catheter & 20$^{*}$ / 10$^{*}$ / 15$^{*}$ & 17 / 9 / 13 \\
RA mean pressure & catheter & 5$^{*}$ & 5 \\
LV / RV systolic pressure & catheter & 94 / 25 & -- \\
LV--RV peak systolic gradient & LV $-$ RV & 69 & -- \\
\midrule
\multicolumn{4}{l}{\emph{Flows and resistance}} \\
Pulmonary flow $Q_p$ (L/min) & Fick & 4.087 & -- \\
$Q_p/Q_s$ (--) & Fick / echocardiography$^{\mathrm{a}}$ & 1.194$^{*}$ & 0.90 \\
Systemic flow $Q_s$ (L/min) & $Q_p/(Q_p/Q_s)$ / echocardiography$^{\mathrm{a}}$ & 3.423$^{*}$ & 2.54 \\
Shunt flow (L/min) & $Q_p - Q_s$ & 0.664$^{*}$ & -- \\
SVR (WU) & $(\overline{P} - P_{\mathrm{RA}})/Q_s$ & 19.38 & 27.56 \\
\midrule
\multicolumn{4}{l}{\emph{Ventricular volumes and function}} \\
LVEDV / LVESV (mL) & echocardiography, biplane$^{\mathrm{a}}$ & -- & 36.13 / 13.80$^{\mathrm{b}}$ \\
RVEDV / RVESV (mL) & echocardiography, 3D & -- & 30.5 / 12.0 \\
LVEF / RVEF (--) & from volumes & -- & 0.618 / 0.607 \\
\bottomrule
\end{tabular}}
\par\vspace{2pt}
\begin{minipage}{0.97\linewidth}
\raggedright\footnotesize
\emph{Abbreviations:} DAo, descending aorta; EDV/ESV, end-diastolic/end-systolic volume; EF, ejection fraction; LV/RV, left/right ventricle; PA, pulmonary artery; RA, right atrium; RFA, right femoral artery; SVR, systemic vascular resistance.\\
\emph{Note:} Values are from the catheterization protocol forms except where marked. $^{\mathrm{a}}$Post-closure value from the medical record. $^{\mathrm{b}}$The protocol form gives 32.0 / 23.6~mL, measured before device release (Section~\ref{sec:data}).
\end{minipage}
\end{table}

% ---------------------------------------------------------------------
\begin{figure}[!htbp]
\centering
\begin{tikzpicture}[
  node distance=4.5mm,
every node/.append style={execute at begin node={\hyphenpenalty=10000 \exhyphenpenalty=10000}},
  box/.style={draw, rounded corners=1pt, align=center, font=\footnotesize, text width=7.0cm, inner sep=3pt, fill=white},
  sbox/.style={draw, dashed, rounded corners=1pt, align=center, font=\footnotesize, text width=4.2cm, inner sep=3pt},
  arr/.style={-{Stealth[length=2mm]}, thick}]
\node[box] (clin) {\textbf{Clinical record} (Table~\ref{tab:patient_data})\\ catheterization and echocardiography};
\node[box, below=of clin] (scale) {\textbf{Weight-based allometric scaling} (Eq.~\ref{eq:scaling})};
\node[box, below=of scale] (seed) {\textbf{Clinical seeding and baseline simulation}\\ (Table~\ref{tab:seeding})};
\node[box, below=of seed] (gsa) {\textbf{Global sensitivity analysis}\\ sparse PCE, $N=128$; retain $S_T \ge 0.10$ (7 of 14)};
\node[box, below=of gsa] (cal) {\textbf{Staged multi-start calibration}\\ 7 parameters, 6 starts, \texttt{fmincon}};
\node[box, below=of cal] (eval) {\textbf{Pre-closure comparison}\\ relative error, $z$ and RMSE};
\node[box, below=of eval] (pred) {\textbf{Post-closure prediction}\\ $A_{\mathrm{VSD}} = 0$, HR = 108 bpm, no refitting};
\node[box, below=of pred, fill=gray!15] (cmp) {\textbf{Comparison with the withheld}\\ \textbf{post-closure record}};
\node[box, below=of cmp] (pcal) {\textbf{Post-closure calibration}\\ 6 parameters, after the comparison};
\foreach \a/\b in {clin/scale, scale/seed, seed/gsa, gsa/cal, cal/eval, eval/pred, pred/cmp, cmp/pcal} \draw[arr] (\a) -- (\b);
\node[sbox, right=8mm of eval] (noise) {\textbf{Measurement-noise simulation}\\ (Section~\ref{sec:noise_methods})\\ $M=20$ perturbed records,\\ each recalibrated and\\ used to predict};
\draw[arr, dashed] ([xshift=-4mm]noise.north) |- (cal.east);
\draw[arr, dashed] ([xshift=-4mm]noise.south) |- (pred.east);
\end{tikzpicture}
\caption{Patient-specific modeling workflow. The clinical record scales and seeds the model, the sensitivity screen selects the seven fitted parameters, staged multi-start optimization calibrates the pre-closure state, and the calibrated parameters are applied unchanged with the defect closed to predict the withheld post-closure state. The measurement-noise simulation repeats calibration and prediction on perturbed records. After the comparison, the model is calibrated to the post-closure record. PCE, polynomial chaos expansion; $S_T$, total-order Sobol' index; HR, heart rate.}
\label{fig:workflow}
\end{figure}


\FloatBarrier
\subsection{Mathematical model}
\label{sec:model}

\subsubsection{Structure}
The circulation is the closed-loop lumped-parameter model of Regazzoni et al.~\cite{Regazzoni2022}, in the version of Tonini et al.\ with oxygenated and non-oxygenated pulmonary capillary branches~\cite{Tonini2024,Tonini2025}, extended with an interventricular shunt by Valenti and Ferrero et al.~\cite{Valenti2023,Ferrero2026Shunts} (Figure~\ref{fig:schematic}). It contains four time-varying elastance chambers, a systemic circuit with arterial, capillary and venous compartments, a pulmonary circuit with an arterial compartment, two parallel capillary branches and a venous compartment, four valves and the shunt. Table~\ref{tab:components} lists the source of each component and what we changed. The model has 14 state variables,
\begin{equation}
\begin{split}\mathbf{X} = \big[ & V_{\mathrm{RA}}, V_{\mathrm{RV}}, V_{\mathrm{LA}}, V_{\mathrm{LV}}, V_{\mathrm{SAR}}, Q_{\mathrm{SAR}}, V_{\mathrm{SC}}, \\
  & V_{\mathrm{SVEN}}, Q_{\mathrm{SVEN}}, V_{\mathrm{PAR}}, Q_{\mathrm{PAR}}, P_{\mathrm{PC}}, V_{\mathrm{PVEN}}, Q_{\mathrm{PVEN}} \big]^{\top},\end{split}
\label{eq:state}
\end{equation}
with volumes $V$ (mL), flows $Q$ (mL/s) and the pulmonary capillary pressure $P_{\mathrm{PC}}$ (mmHg). The chamber volumes obey
\begin{equation}
\begin{aligned}
\dot{V}_{\mathrm{RA}} &= Q_{\mathrm{SVEN}} - Q_{\mathrm{TV}}, &
\dot{V}_{\mathrm{RV}} &= Q_{\mathrm{TV}} + Q_{\mathrm{VSD}} - Q_{\mathrm{PV}}, \\
\dot{V}_{\mathrm{LA}} &= Q_{\mathrm{PVEN}} - Q_{\mathrm{MV}}, &
\dot{V}_{\mathrm{LV}} &= Q_{\mathrm{MV}} - Q_{\mathrm{AV}} - Q_{\mathrm{VSD}},
\end{aligned}
\label{eq:chambers}
\end{equation}
where $Q_{\mathrm{TV}}$, $Q_{\mathrm{PV}}$, $Q_{\mathrm{MV}}$ and $Q_{\mathrm{AV}}$ are the flows through the tricuspid, pulmonary, mitral and aortic valves and $Q_{\mathrm{VSD}}$ is the shunt flow, positive from left to right. The vascular equations are given in Appendix~\ref{app:circuit} and the nominal adult values in Appendix~\ref{app:nominal_parameters}. The nominal values were obtained by manual tuning, starting from the adult values of Valenti~\cite{Valenti2023}, Zhang et al.~\cite{zhang2019multiscale} and Bozkurt~\cite{Bozkurt2019}, until the weight-scaled model of a healthy 3-year-old (14~kg, 98~cm, the WHO median for age~\cite{WHOgrowth}) matched pediatric reference ranges for heart rate~\cite{Joyce2025}, systemic blood pressure~\cite{Flynn2017} and pulmonary arterial and right atrial pressure~\cite{MalakanRad2024}. Supplementary Section~S6 lists every nominal value beside its three sources and every output of the healthy-child model beside its reference range. The external pressure on the chambers is zero, the two ventricles are not mechanically coupled through the septum, and valve regurgitation is not modeled.

\begin{table}[!htbp]
\centering\small
\caption{Model components, their source and the changes made in this work.}
\label{tab:components}
\resizebox{\linewidth}{!}{\begin{tabular}{lll}
\toprule
\textbf{Component} & \textbf{Source} & \textbf{Change in this work} \\
\midrule
Chambers and activation (Eqs.~\ref{eq:elastance}--\ref{eq:activation}) & \cite{Regazzoni2022,Valenti2023} & none \\
Arterial and venous compartments (Appendix~\ref{app:circuit}) & \cite{Tonini2024,Tonini2025} & none \\
Systemic capillary compartment & \cite{Tonini2025} & resistance split into two halves around $C_{\mathrm{SC}}$ \\
Pulmonary capillary compartment, two branches & \cite{Tonini2024,Tonini2025} & none \\
Valves (Eq.~\ref{eq:valve}) & smooth valve resistance~\cite{Regazzoni2022,Tonini2025} & tanh-gated conductance \\
Interventricular shunt (Eq.~\ref{eq:orifice}) & linear resistance~\cite{Valenti2023,Ferrero2026Shunts} & orifice law with $C_d$ and measured diameter \\
Capillary and venous compliances (Eq.~\ref{eq:rc}) & constant & follow their resistance during calibration \\
Atrial and venous pressures & -- & bounded below at $-5$~mmHg \\
Nominal parameter values (Appendix~\ref{app:nominal_parameters}) & \cite{Valenti2023,zhang2019multiscale,Bozkurt2019} & tuned to a healthy 3-year-old (Supplementary Section~S6) \\
\bottomrule
\end{tabular}}
\end{table}

\begin{figure}[!htbp]
\centering
\includegraphics[width=\linewidth]{model_schematic.pdf}
\caption{Electrical analogue of the closed-loop model. Chambers are time-varying elastances, valves are smooth diodes, and the interventricular defect is a nonlinear orifice; $Q_{\mathrm{VSD}}$ is positive from LV to RV. $C_{\mathrm{PC}} = C_{\mathrm{PCOX}} + C_{\mathrm{PCNO}}$. Blue elements are calibrated, directly or through a resistance scale; $C_{\mathrm{SC}}$, $C_{\mathrm{SVEN}}$, $C_{\mathrm{PVEN}}$ and the oxygenated part of $C_{\mathrm{PC}}$ follow their resistance (Eq.~\ref{eq:rc}); all other elements are held at their allometrically scaled values. MV, AV, TV, PV: mitral, aortic, tricuspid and pulmonary valves.}
\label{fig:schematic}
\end{figure}

\FloatBarrier
\subsubsection{Cardiac chambers}
Chamber pressure follows the time-varying elastance model~\cite{Suga1974,Regazzoni2022},
\begin{equation}
P_{ch}(t) = E_{ch}(t)\left(V_{ch}(t) - V_{0,ch}\right), \qquad
E_{ch}(t) = E_{\mathrm{B},ch} + E_{\mathrm{A},ch}\, e_{ch}(t),
\label{eq:elastance}
\end{equation}
with unstressed volume $V_{0,ch}$, passive elastance $E_{\mathrm{B},ch}$, active elastance $E_{\mathrm{A},ch}$ and activation $e_{ch}(t)\in[0,1]$. Ventricular activation over the cardiac phase $\phi = \mathrm{mod}(t, T_{\mathrm{HB}})$, with $T_{\mathrm{HB}} = 60/\mathrm{HR}$, is the piecewise cosine of Regazzoni et al.~\cite{Regazzoni2022},
\begin{equation}
e_v(\phi) =
\begin{cases}
\tfrac{1}{2}\left[1-\cos\left(\pi\phi/T_{vc}\right)\right], & 0 \le \phi \le T_{vc},\\[3pt]
\tfrac{1}{2}\left[1+\cos\left(\pi(\phi-T_{vc})/T_{vr}\right)\right], & T_{vc} < \phi \le T_{vc}+T_{vr},\\[3pt]
0, & \text{otherwise},
\end{cases}
\label{eq:activation}
\end{equation}
where $T_{vc}$ and $T_{vr}$ are the contraction and relaxation durations. Atrial activation uses the same form with $\tau = \mathrm{mod}(\phi - t_{ac}, T_{\mathrm{HB}})$ in place of $\phi$ and the atrial durations $T_{ac}$ and $T_{ar}$, where $t_{ac}$ is the onset of atrial contraction. All timings are fixed fractions of $T_{\mathrm{HB}}$ (Appendix~\ref{app:nominal_parameters}), so they rescale with the measured heart rate of each state.

\FloatBarrier
\subsubsection{Valves and interventricular shunt}
\label{sec:valves}
Regazzoni et al.\ and Tonini et al.\ model each valve as a resistance that changes smoothly between an open and a closed value~\cite{Regazzoni2022,Tonini2025}. We use a tanh-gated conductance instead,
\begin{equation}
g = \tfrac{1}{2} + \tfrac{1}{2}\tanh\!\left(\frac{\Delta P}{\varepsilon_{\mathrm{v}}}\right), \qquad
Q_{\mathrm{valve}} = \Delta P\left(\frac{g}{R_{\mathrm{open}}} + \frac{1-g}{R_{\mathrm{closed}}}\right),
\label{eq:valve}
\end{equation}
with transvalvular gradient $\Delta P$ and $\varepsilon_{\mathrm{v}} = 0.5$~mmHg. The nominal resistances are $R_{\mathrm{open}} = 4.0\times10^{-3}$ and $R_{\mathrm{closed}} = 9.42\times10^{4}$~mmHg\,s/mL; after scaling (Section~\ref{sec:scaling}), $R_{\mathrm{open}} = 9.14\times10^{-3}$~mmHg\,s/mL for P1.

Valenti and Ferrero et al.\ represent the defect as a linear resistance~\cite{Valenti2023,Ferrero2026Shunts}. We replace it with the Gorlin orifice relation~\cite{Gorlin1951}, so that the shunt depends on the measured defect area, and add a sign-preserving factor that makes the relation differentiable at zero gradient,
\begin{equation}
Q_{\mathrm{VSD}} = 10^{6}\, C_d\, A_{\mathrm{VSD}} \sqrt{\frac{2\,\gamma\,\lvert\Delta P_{\mathrm{v}}\rvert}{\rho}}\;
\frac{\Delta P_{\mathrm{v}}}{\sqrt{\Delta P_{\mathrm{v}}^{2} + \varepsilon_{\mathrm{s}}^{2}}},
\qquad A_{\mathrm{VSD}} = \pi\left(\frac{D_{\mathrm{VSD}}}{2}\right)^{2},
\label{eq:orifice}
\end{equation}
where $Q_{\mathrm{VSD}}$ is in mL/s, $A_{\mathrm{VSD}}$ in m$^2$ (the factor $10^{6}$ converts m$^3$/s to mL/s), $\Delta P_{\mathrm{v}} = P_{\mathrm{LV}} - P_{\mathrm{RV}}$ in mmHg, $\gamma = 133.322$~Pa/mmHg, $\rho = 1060$~kg/m$^3$ and $\varepsilon_{\mathrm{s}} = 0.1$~mmHg. The discharge coefficient $C_d$ lumps the defect shape and the contraction of the jet; it is the only shunt parameter that is calibrated. Closure sets $A_{\mathrm{VSD}} = 0$, which gives $Q_{\mathrm{VSD}} = 0$ exactly.

\subsubsection{Coupling of compliance to resistance}
\label{sec:rc}
A change in the resistance of a capillary or venous compartment represents a change in vessel radius, which also changes the compliance of that compartment. During calibration we therefore recompute these compliances from their resistances with the relation of Corsini et al.~\cite{Corsini2014}, which follows from the dependence of vascular resistance and compliance on vessel radius and length~\cite{Pennati2000},
\begin{equation}
\frac{C_i}{C_i^{\mathrm{s}}} = \left(\frac{R_i}{R_i^{\mathrm{s}}}\right)^{-4/3}, \qquad i \in \{\mathrm{SC},\,\mathrm{SVEN},\,\mathrm{PCOX},\,\mathrm{PVEN}\},
\label{eq:rc}
\end{equation}
where $C_i^{\mathrm{s}}$ and $R_i^{\mathrm{s}}$ are the allometrically scaled values. The ratio $C_i/C_i^{\mathrm{s}}$ is limited to 0.05--8. The arterial compliances $C_{\mathrm{SAR}}$ and $C_{\mathrm{PAR}}$ are calibrated directly and are not coupled. The sensitivity screen does not use Eq.~(\ref{eq:rc}): it samples every resistance and compliance independently (Section~\ref{sec:gsa}).

\FloatBarrier
\subsection{Pediatric scaling and initialization}
\label{sec:scaling}

\subsubsection{Allometric scaling}
Each adult reference parameter was scaled to the weight of P1 with the allometric law of Zhang et al.~\cite{zhang2019multiscale},
\begin{equation}
\theta_i = \theta_{i,\mathrm{ref}} \left(\frac{W_{\mathrm{p}}}{W_{\mathrm{ref}}}\right)^{\alpha_i},
\label{eq:scaling}
\end{equation}
where $W_{\mathrm{ref}} = 70$~kg is the adult reference weight, $W_{\mathrm{p}} = 13.4$~kg the weight of P1, and $\alpha_i$ an exponent that depends on the parameter type (Table~\ref{tab:exponents}). Zhang et al.\ give no exponent for inertance; we used $\alpha = -1$ from Lundquist et al.~\cite{Lundquist2025}. Heart rate was not scaled but set to the measured value of each state.

\begin{table}[!htbp]
\centering\small
\caption{Weight-based allometric exponents used in Eq.~(\ref{eq:scaling}).}
\label{tab:exponents}
\begin{tabular}{llc}
\toprule
\textbf{Class} & \textbf{Parameters} & $\alpha_i$ \\
\midrule
Left-heart elastance & $E_{\mathrm{A/B,LV}}$, $E_{\mathrm{A/B,LA}}$ & $-0.50$ \\
Right-heart elastance & $E_{\mathrm{A/B,RV}}$, $E_{\mathrm{A/B,RA}}$ & $-0.75$ \\
Systemic resistance & $R_{\mathrm{SAR}}$, $R_{\mathrm{SC}}$, $R_{\mathrm{SVEN}}$ & $-0.475$ \\
Pulmonary resistance & $R_{\mathrm{PAR}}$, $R_{\mathrm{PCOX}}$, $R_{\mathrm{PCNO}}$, $R_{\mathrm{PVEN}}$ & $-0.70$ \\
Compliance & all vascular $C$ & $+1.00$ \\
Unstressed volume & all $V_0$ & $+0.80$ \\
Inertance & all $L$ & $-1.00$\,$^{\mathrm{a}}$ \\
Valve open resistance & $R_{\mathrm{open}}$ & $-0.50$ \\
\bottomrule
\end{tabular}
\par\vspace{2pt}
\begin{minipage}{0.8\linewidth}
\raggedright\footnotesize
\emph{Note:} All exponents are from Zhang et al.~\cite{zhang2019multiscale} except $^{\mathrm{a}}$, from Lundquist et al.~\cite{Lundquist2025}. Heart rate is not scaled (set to the measured value of each state).
\end{minipage}
\end{table}

\subsubsection{Initial conditions and clinically seeded parameters}
Table~\ref{tab:seeding} lists every initial condition and every parameter seeded from the clinical record, with its formula and value for P1. The total blood volume is $\mathrm{BV} = 70$~mL/kg~$\times\,W_{\mathrm{p}} = 938$~mL, the standard estimate for children older than three years~\cite{Cote2024}. The systemic venous unstressed volume closes the volume budget,
\begin{equation}
V_{0,\mathrm{SVEN}} = \mathrm{BV} - \sum_{i \neq \mathrm{SVEN}} V_i(0) - P_{\mathrm{nom,SVEN}}\,C_{\mathrm{SVEN}},
\label{eq:sven}
\end{equation}
where the sum runs over the four chambers and the other eight vascular volume states, bounded to $[\max(0.01\,\mathrm{BV},\,1~\mathrm{mL}),\ 0.85\,\mathrm{BV}]$. The vascular resistances start at their scaled values.

\begin{table}[!htbp]
\centering\small
\setlength{\tabcolsep}{4pt}
\caption{Initial conditions and clinically seeded parameters for P1. $P_{\mathrm{nom}}$ is the nominal compartment pressure used to place the initial volume.}
\label{tab:seeding}
\resizebox{\linewidth}{!}{\begin{tabular}{llll}
\toprule
\textbf{Quantity} & \textbf{Definition} & \textbf{Formula} & \textbf{Value} \\
\midrule
\multicolumn{4}{l}{\emph{Initial conditions}} \\
BV & total blood volume & $70$~mL/kg $\times\,W_{\mathrm{p}}$ & 938 mL \\
$V_i(0)$, vascular & initial vascular volume & $V_{0,i} + P_{\mathrm{nom},i}\,C_i$ & -- \\
$V_{ch}(0)$ & initial chamber volume & $V_{0,ch} + P_{\mathrm{ED},ch}/E_{\mathrm{B},ch}$ & -- \\
$P_{\mathrm{nom,SAR}}$ & systemic arterial & mean arterial pressure & 71.33 mmHg \\
$P_{\mathrm{nom,SC}}$ & systemic capillary & $\max(0.5\,P_{\mathrm{nom,SAR}},\,15)$ & 35.7 mmHg \\
$P_{\mathrm{nom,SVEN}}$, $P_{\mathrm{ED,RA}}$, $P_{\mathrm{ED,RV}}$ & venous and right-heart filling & mean RA pressure & 5 mmHg \\
$P_{\mathrm{nom,PAR}}$ & pulmonary arterial & mean PA pressure & 15 mmHg \\
$P_{\mathrm{nom,PVEN}}$, $P_{\mathrm{ED,LA}}$, $P_{\mathrm{ED,LV}}$ & left-heart filling & assumed (LA pressure not measured) & 6 / 6 / 8 mmHg \\
$P_{\mathrm{PC}}(0)$ & pulmonary capillary & $(P_{\mathrm{nom,PAR}}+P_{\mathrm{nom,PVEN}})/2$ & 10.5 mmHg \\
\midrule
\multicolumn{4}{l}{\emph{Clinically seeded parameters}} \\
$C_{\mathrm{SAR}}$ & systemic arterial compliance & $(Q_s/\mathrm{HR})/(P_{\mathrm{sys}}-P_{\mathrm{dia}})_{\mathrm{RFA}}$ & 0.669 mL/mmHg \\
$C_{\mathrm{PAR}}$ & pulmonary arterial compliance & $(Q_p/\mathrm{HR})/(P_{\mathrm{sys}}-P_{\mathrm{dia}})_{\mathrm{PA}}$ & 3.434 mL/mmHg \\
$C_d$ & discharge coefficient & Eq.~(\ref{eq:orifice}) solved for $C_d$ at $Q_{\mathrm{shunt}}$ and $\Delta P_{\mathrm{v}} = \tfrac12 \Delta P_{\mathrm{LV-RV,peak}}$ & 0.356 \\
\bottomrule
\end{tabular}}
\par\vspace{2pt}
\begin{minipage}{0.97\linewidth}
\raggedright\footnotesize
\emph{Note:} $C_{\mathrm{SAR}}$ and $C_{\mathrm{PAR}}$ are set to stroke volume over pulse pressure, following Colebank et al.~\cite{Colebank2026}. The $C_d$ seed uses half of the measured peak LV--RV systolic gradient (69~mmHg) as a representative gradient over the cycle.
\end{minipage}
\end{table}

\FloatBarrier
\subsection{Numerical solution}
\label{sec:numerics}
The system was integrated with \texttt{ode15s}~\cite{Shampine1997} in MATLAB R2025a, with relative tolerance $10^{-5}$, absolute tolerance $10^{-6}$, maximum step $T_{\mathrm{HB}}/20$ and the sparsity pattern of the Jacobian supplied. The solver advanced the model in batches of five cardiac cycles and compared the peak value of each state between consecutive batches. Periodic steady state was declared when every volume state changed by less than 0.1~mL, $P_{\mathrm{PC}}$ by less than 0.1~mmHg and every flow state by less than 0.5\% of its peak (0.1~mL/s when the peak was below 0.1~mL/s), within at most 80 cycles. Outputs were computed from the last two cycles. During calibration, each simulation was limited to 30--60 cycles with a steady-state tolerance of 0.5~mL and 0.5~mmHg to reduce cost, and a simulation that did not reach steady state added a penalty of $10^{3}$ to the objective. The reported states were recomputed with the full settings and reached steady state. The atrial and venous pressures are bounded below at $-5$~mmHg to avoid solver failure from non-physiological initial states; the bound is not active in the reported states, where the lowest atrial pressure is 2.9~mmHg and the lowest venous pressure 8.5~mmHg. Tightening the tolerances tenfold and a hundredfold changed no output by more than 0.07\%, reducing the maximum step to $T_{\mathrm{HB}}/50$ changed no output, running 20 further cycles after steady state changed no output by more than 0.05\%, and total blood volume was conserved over the cycle to within $10^{-4}$~mL (Supplementary Section~S4).

\FloatBarrier
\subsection{Global sensitivity analysis}
\label{sec:gsa}
Variance-based global sensitivity analysis~\cite{Sobol2001,Saltelli2010} ranked the parameters by their influence on the model outputs, following recommendations for cardiovascular models~\cite{Eck2016} and the approach of Tonini et al.\ for this model~\cite{Tonini2025}. We fitted sparse polynomial chaos expansions (PCE)~\cite{Blatman2011} in UQLab~2.2~\cite{Marelli2014} (least-angle regression, degree 1--3, hyperbolic truncation $q = 0.75$) to one Halton design of $N = 128$ points over $d = 21$ independent uniform inputs. The input ranges were $[0.4,\,2.5]\,x_0$ for resistances, $[0.5,\,2.0]\,x_0$ for compliances and elastances, $[0.6,\,1.7]\,x_0$ for unstressed volumes and $[0.8,\,1.2]\,x_0$ for $C_d$, where $x_0$ is the scaled and seeded value. The total-order Sobol' index of each input,
\begin{equation}
S_{T_i} = 1 - \frac{\mathrm{Var}\!\left(\mathbb{E}\left[Y \mid \mathbf{X}_{\sim i}\right]\right)}{\mathrm{Var}(Y)},
\label{eq:sobol}
\end{equation}
was computed analytically from the expansion coefficients~\cite{Sudret2008} for 25 outputs. To keep the design affordable, each training simulation ran for at most 10 cycles with a steady-state tolerance of 1.0~mmHg and 1.0~mL, and outputs were taken from cycles 9--10.

The 21 inputs contain the 14 calibration candidates. The systemic resistance scale multiplies $R_{\mathrm{SAR}}$ and $R_{\mathrm{SC}}$, and the pulmonary resistance scale multiplies $R_{\mathrm{PAR}}$, $R_{\mathrm{PCOX}}$ and $R_{\mathrm{PVEN}}$; each scale takes the largest index of its members. The other four inputs, $C_{\mathrm{SVEN}}$, $C_{\mathrm{PVEN}}$, $V_{0,\mathrm{LA}}$ and $V_{0,\mathrm{RA}}$, are not calibration candidates: the two venous compliances follow their resistances during calibration (Eq.~\ref{eq:rc}), and the atrial unstressed volumes were held at their scaled values. A candidate was retained if $S_{T} \ge 0.10$ for at least one of the nine calibration targets (Section~\ref{sec:comparison}), the threshold Valenti used to select calibration parameters in patients with septal defects~\cite{Valenti2023}. The screen fixes the calibrated parameter set; no later stage adds a parameter. The screen was repeated at the calibrated operating point with the same relative input ranges.

\FloatBarrier
\subsection{Calibration}
\label{sec:calibration}

\subsubsection{Objective function}
Calibration minimized an objective with two parts: a fit term that measures how far the model is from the nine calibration targets, and penalty terms that keep the solution physiological and the parameters plausible,
\begin{equation}
J(\mathbf{x}) = \underbrace{\sum_{k} w_k\left(\frac{\varepsilon_k}{s_k}\right)^{2} + J_{\mathrm{sys}}}_{\text{fit to the targets}}
\; + \; \underbrace{J_{\mathrm{shape}} + J_{\mathrm{shunt}} + J_{\mathrm{gate}} + J_{\mathrm{guard}}}_{\text{penalties on large target errors}}
\; + \; \underbrace{J_{\mathrm{prior}} + J_{\mathrm{bound}} + J_{\mathrm{range}}}_{\text{penalties on implausible solutions}},
\label{eq:objective}
\end{equation}
where $\varepsilon_k = (y_k - y_k^{\mathrm{clin}})/y_k^{\mathrm{clin}}$ is the relative error of target $k$. Relative errors are used so that pressures and flows, which have different units and magnitudes, can be added in one sum. Each error is scaled by $s_k = 0.05$ for mean PA pressure and $Q_p/Q_s$ and $s_k = 0.10$ for the systolic and diastolic pressures and shunt flow, and weighted by $w_k$ (Supplementary Table~S3).

Mean RA pressure, mean arterial pressure and $Q_s$ describe the same systemic load, so they enter one combined term, $J_{\mathrm{sys}}$, scaled by their measurement uncertainty, so that a single systemic mismatch is not counted three times. The remaining terms do not fit new data. Four of them add extra cost when a target error becomes large: on the systemic and pulmonary pulse pressure ($J_{\mathrm{shape}}$), on the agreement between the orifice flow and $Q_p - Q_s$ ($J_{\mathrm{shunt}}$), and on any target with a relative error above 10\% ($J_{\mathrm{gate}}$) or a pressure error above 12\% ($J_{\mathrm{guard}}$). Three of them keep the solution plausible: $J_{\mathrm{prior}}$ penalizes drift of a parameter from its scaled value, $J_{\mathrm{bound}}$ penalizes parameters that approach a bound, and $J_{\mathrm{range}}$ penalizes outputs without a clinical value, such as chamber volumes, that leave pediatric physiological ranges, and simulations that do not reach steady state. The formula and weight of every term are given in Supplementary Section~S3.

\subsubsection{Bounds}
Each calibrated parameter has a lower and an upper bound (Table~\ref{tab:bounds}): $[0.25,\,2.8]$ for the systemic and $[0.45,\,2.8]$ for the pulmonary resistance scale; $[0.40,\,2.8]$ and $[0.55,\,2.6]$ times the scaled value for $R_{\mathrm{SVEN}}$ and $E_{\mathrm{B,RV}}$; $[0.75,\,1.35]$ and $[0.70,\,1.45]$ times the clinical seed for $C_{\mathrm{SAR}}$ and $C_{\mathrm{PAR}}$, with the lower bound extended to just below the scaled value (the scaled value minus 5\% of it or 2\% of the interval, whichever is larger) so that the scaled value stays feasible; and $[0.8\times$ seed, $1.05\times$ adult reference$]$ for $C_d$. Each interval is centered on a value with a published basis: the allometrically scaled value~\cite{zhang2019multiscale} for the resistance scales, $R_{\mathrm{SVEN}}$ and $E_{\mathrm{B,RV}}$; the clinical seed, stroke volume over pulse pressure~\cite{Colebank2026}, for $C_{\mathrm{SAR}}$ and $C_{\mathrm{PAR}}$; and the adult reference value~\cite{Valenti2023} for the upper bound of $C_d$. The widths of the intervals were chosen to cover the physiological range of each parameter class.

\subsubsection{Staged multi-start optimization}
For each starting point, six bound-constrained subproblems were solved in sequence with \texttt{fmincon} (Table~\ref{tab:stages}). Stages A--E used the interior-point algorithm with a limited-memory BFGS Hessian, forward finite differences with step $10^{-5}$, optimality tolerance $10^{-5}$, step tolerance $10^{-6}$, and at most 300 function evaluations and 40 iterations. Stage F used sequential quadratic programming on a separate score that penalizes each calibration target with a relative error above 10\% (forward finite differences with step $10^{-4}$, at most 900 function evaluations and 120 iterations, two starts: the current point and the clinical seed). Stages D, E and F were kept only if they passed the acceptance rule in Table~\ref{tab:stages}.

\begin{table}[!htbp]
\centering\small
\setlength{\tabcolsep}{3pt}
\caption{Calibration stages for each starting point.}
\label{tab:stages}
\resizebox{\linewidth}{!}{\begin{tabular}{llll}
\toprule
\textbf{Stage} & \textbf{Free parameters} & \textbf{Targets} & \textbf{Kept if} \\
\midrule
A & $R_{\mathrm{sys}}$, $R_{\mathrm{SVEN}}$, $R_{\mathrm{pul}}$, $C_{\mathrm{SAR}}$, $C_{\mathrm{PAR}}$, $C_d$ & RA, RFA and mean PA pressure, $Q_p/Q_s$, $Q_s$, shunt flow & always \\
B & $E_{\mathrm{B,RV}}$ & mean RA pressure & always \\
C & all seven & all calibration targets & always \\
D & $R_{\mathrm{sys}}$, $R_{\mathrm{SVEN}}$, $C_{\mathrm{SAR}}$, $E_{\mathrm{B,RV}}$ & mean arterial, RA and PA pressure, $Q_s$, $Q_p/Q_s$ & RMSE and number of mean-pressure and flow targets above 10\% do not increase \\
E & all seven, $\lambda_{\mathrm{p}} = 2.25$, $\lambda_{\mathrm{b}} = 120$ & all calibration targets & fewer plausibility warnings; RMSE rises by $< 0.005$ \\
F & all seven (SQP, gate score) & all calibration targets & RMSE or number of targets above 10\% falls \\
\bottomrule
\end{tabular}}
\end{table}

Six starting points were used: the clinical seed and five points of a scrambled Sobol' sequence (Matou\v{s}ek affine--Owen scrambling, first point skipped) mapped uniformly into the bounds, generated with the seed 101, fixed in advance. One result was selected in three ordered steps: fewest of the five mean-pressure and flow targets with a relative error above 10\%, then fewest of the nine calibration targets above 10\%, then lowest RMSE. The 10\% value in the objective and in this rule is a setting of the optimizer only; the Results report the relative errors without a threshold. Repeating the pipeline with the seeds 202 and 303 returned the same parameter set (Supplementary Section~S1).

\FloatBarrier
\subsection{Post-closure prediction}
\label{sec:post_methods}
The post-closure state is a prediction, not a fit. The calibrated pre-closure parameters, including the compliances set by Eq.~(\ref{eq:rc}), were applied unchanged with the defect closed ($A_{\mathrm{VSD}} = 0$) and the heart rate set to its post-closure value (108~bpm), which rescales all chamber timings. No parameter was refitted and no post-closure measurement entered calibration. The prediction therefore represents the circulation immediately after closure, with the chamber properties of the pre-closure state.

\subsection{Post-closure calibration}
\label{sec:post_calib_methods}
After the prediction had been compared with the post-closure record, the model was calibrated a second time, to the post-closure measurements. This calibration does not change the prediction; it estimates the parameters of the circulation after closure so that they can be compared with the pre-closure parameters. The procedure of Section~\ref{sec:calibration} was repeated with the defect closed and the heart rate set to 108~bpm: the same scaling, clinical seeding from the post-closure record, bound widths, weights, six starting points and staged optimization. The free parameters were the six pre-closure parameters other than $C_d$, which has no effect once the defect is closed; the sensitivity screen was not repeated. The targets were the seven post-closure catheter pressures and $Q_s$. $Q_p/Q_s$ was not a target, because a value below one cannot be reached without a shunt, and the shunt terms were removed from the objective. The ventricular volumes were compared but not fitted, as before closure.

\FloatBarrier
\subsection{Comparison between model and measurements}
\label{sec:comparison}
\label{sec:statistics}
Each measured quantity is reported with its measured value, model value and relative error $\varepsilon_k$, in the same order before and after closure. The nine calibration targets are summarized by the root-mean-square of their relative errors,
\begin{equation}
\mathrm{RMSE} = \sqrt{\frac{1}{9}\sum_{k=1}^{9}\varepsilon_k^{2}},
\label{eq:rmse}
\end{equation}
which excludes shunt flow because it is the difference of two calibrated flows. No acceptance threshold is applied.

Where an uncertainty $\sigma_k$ is declared (column $\sigma$ in Tables~\ref{tab:pre_targets} and~\ref{tab:post}), the standardized residual $z_k = (y_k^{\mathrm{model}} - y_k^{\mathrm{clin}})/\sigma_k$ is also given. The values of $\sigma_k$ follow the spread of repeated catheter readings for RA and PA pressures (5\%), the difference between the RFA and descending-aortic readings for systolic and diastolic systemic pressure (10\%) and the error of Fick output in children for pre-closure $Q_s$ (14.6\%)~\cite{Ralston2023}. For the remaining quantities we assumed 5\% for mean arterial pressure, $Q_p/Q_s$ and post-closure $Q_s$, 10\% for ventricular volumes and ejection fractions, and 20\% for the derived shunt flow and SVR.

Pressure--volume loops were characterized by end-diastolic and end-systolic volume (maximum and minimum volume), stroke volume, ejection fraction, end-diastolic pressure, end-systolic pressure, peak pressure and stroke work $\mathrm{SW} = \oint P\,dV$. End-systolic pressure is taken at the maximum of $P/V$, because with an open defect the LV keeps emptying into the RV after pressure has fallen. Practical identifiability is reported in Supplementary Section~S2.

\FloatBarrier
\subsection{Measurement-noise simulation}
\label{sec:noise_methods}
We propagated measurement uncertainty to the calibrated parameters and the post-closure prediction by Monte Carlo simulation. In each of $M = 20$ replicates, Gaussian noise with standard deviation $\sigma_k$ (column $\sigma$ in Table~\ref{tab:pre_targets}) was added independently to every measured pre-closure calibration target; mean arterial pressure and shunt flow were recomputed from the noisy values, and physically impossible draws (e.g.\ $Q_p/Q_s \le 1$) were redrawn. $M = 20$ was set by computational cost, about 85~min per recalibration.

Each noisy record was recalibrated with the procedure of Section~\ref{sec:calibration}, with the parameter set of the original calibration, and used to predict the post-closure state, which was compared with the original post-closure record. A replicate was excluded when its seeded $C_d$ fell outside the $C_d$ bounds, so that the calibration could not start, or when its calibration ended with a different set of free parameters. For each parameter and predicted quantity we report the median, coefficient of variation and range across the remaining replicates.

\FloatBarrier
\subsection{Implementation}
The model, calibration and analyses were implemented in MATLAB R2025a (MathWorks, Natick, MA, USA) with the Optimization and the Statistics and Machine Learning Toolboxes and UQLab~2.2~\cite{Marelli2014}; figures were drawn in R~4.5.2 with ggplot2. One calibration took about 85~min on a laptop with an AMD Ryzen~5 5600H processor and 16~GB of RAM. One script (\texttt{run\_all.m}) reproduces the pre-closure pipeline, the post-closure prediction and every exported table, a second (\texttt{run\_noise\_simulation.m}) reproduces the measurement-noise simulation, and a third (\texttt{run\_post\_closure\_calibration.m}) reproduces the post-closure calibration.
